\documentclass{situationist-preprint}

\title{Emotions Represent Values}
\msid{MND-17-a}
\msfamily{MND-17}
\msdiscipline{Philosophy of Mind}
\mstype{Principal manuscript}
\mskeywords{emotion; intentionality; evaluative content; formal objects; fittingness; recalcitrant emotions; perceptual theory}

\begin{document}
\maketitle

\begin{abstract}
Do emotions have content, and if so, what content? We argue that emotions represent evaluative properties: fear represents its object as dangerous, disgust represents its object as contaminating, and each emotion type represents its object as having the evaluative property that is its formal object. The argument starts from the fact that emotions are assessable as fitting or unfitting, and that the fittingness of an instance of a type depends on whether its object has a property specific to the type. We show that these correctness conditions are not those of an accompanying belief, since recalcitrant emotional states conflict with the subject's beliefs, and not those of bodily feelings, which have no correctness conditions concerning the emotion's object. A state whose correctness depends in this way on its object's having a property represents its object as having it. The account is consistent with the view that fear proper also requires the belief that the danger is actual, and with the existence of quasi-emotions that lack such beliefs. We reply to the attitudinal theory, to the problem of objectless moods and to the charge of over-intellectualization.
\end{abstract}

\section{The problem}

Emotions are directed at things: one fears the dog, is disgusted by the spoiled milk, is angry at the colleague. \citet{James1884} identified an emotion with the feeling of the bodily changes that follow the perception of an exciting fact, which makes the directedness of emotion hard to explain, since a feeling of a racing heart is not about a dog. The cognitivist tradition went to the other extreme and identified emotions with evaluative judgments: to fear the dog is to judge it dangerous \citep{Solomon1976,Nussbaum2001}. This explains directedness but not the fact that one can fear what one judges harmless.

Between these, two views have dominated recent work. On the \emph{perceptual} view, emotions are perceptions, or perception-like experiences, of evaluative properties \citep{deSousa1987,Tappolet2016}; on the \emph{embodied appraisal} view, they are perceptions of bodily changes that, in virtue of their function, represent core relational themes such as danger or loss \citep{Prinz2004,Lazarus1991}. On the \emph{attitudinal} view, emotions do not represent evaluative properties at all; they are distinctive attitudes, felt bodily stances toward an object, whose correctness conditions involve evaluative properties without those properties figuring in their content \citep{DeonnaTeroni2012,DeonnaTeroni2015}.

We defend the thesis common to the perceptual and embodied appraisal views: emotions represent evaluative properties. Our argument does not depend on the claim that emotions are perceptions, which we leave open. It depends on the structure of their correctness conditions.

\section{Framework}\label{s:framework}

\begin{definition}[Evaluative property]\label{d:eval}
An \emph{evaluative property} is a property whose possession by an object consists in, or grounds, there being reason for some response to it: being dangerous, contaminating, offensive, a loss, an infringement of a claim.
\end{definition}

\begin{definition}[Formal object]\label{d:formal}
The \emph{formal object} of an emotion type is the property that an instance of the type must take its particular object to have in order to be an instance of that type \citep{Kenny1963}.
\end{definition}

\begin{principle}[Fittingness]\label{p:fit}
Emotions are assessable as fitting or unfitting, and whether an instance of an emotion type is fitting depends on whether its object has a property specific to the type.
\end{principle}

\begin{principle}[Independence from belief]\label{p:belief}
The fittingness of an emotion is not the truth of any belief that accompanies it.
\end{principle}

\begin{principle}[Independence from feeling]\label{p:feel}
The bodily feelings that accompany an emotion have, as such, no correctness conditions concerning the emotion's object.
\end{principle}

\begin{assumption}[Correctness and content]\label{a:content}
If a mental state is correct or incorrect according as its object has or lacks a property $F$, and this is not due to an accompanying belief or to a feeling, then the state itself represents its object as $F$.
\end{assumption}

Principle~\ref{p:fit} records a feature of our practice of assessing emotions which is distinct from moral and prudential assessment. Fear of a harmless spider is unfitting, even if it is prudent to feel it; amusement at a cruel joke may be fitting, in that the joke is funny, even if it is wrong to feel it \citep{DArmsJacobson2000}. The property on which fittingness depends is specific to the type: fear is fitting when its object is dangerous, disgust when its object is contaminating or would contaminate, anger when its object is an offence, sadness when it is a loss.

Principle~\ref{p:belief} is supported by recalcitrant emotional states \citep{Brady2009}. A subject who has studied the safety record of air travel and believes that her flight is safe may nonetheless have, during take-off, the racing heart, the narrowed attention and the felt dread that are fear's, directed at the flight. Her state is unfitting, since the flight is not dangerous; but the belief she holds, that it is safe, is true. If the fittingness of the state were the truth of a belief, the only candidate would be the belief that the flight is dangerous, and she does not have it. AES-05-a holds that fear proper requires the belief that its object exists and threatens, and the recalcitrant state is then not fear proper; it has the affective component of fear without the belief. The argument requires only that this component be assessable as fitting or unfitting, which it is: we say that she should not be frightened, and mean that the flight does not warrant it.

Principle~\ref{p:feel} is the point that defeated James's theory. A feeling of a racing heart is accurate or inaccurate, if at all, about the heart. It is not made correct by the dog's being dangerous, and the same feeling occurs in fear, excitement and the effects of a stimulant.

Assumption~\ref{a:content} links correctness to representation. It is the principle by which we ascribe content to perceptual experiences: an experience is accurate according as the scene is as it is presented, and we infer that it presents the scene as being that way. A state that is correct according as its object has a property, and not because some other state of the subject represents the object, is a state that represents its object as having it.

\section{The argument}\label{s:argument}

\begin{theorem}\label{t:main}
Each emotion type represents its object as having an evaluative property, namely its formal object.
\end{theorem}

\begin{proof}
Let $E$ be an emotion type and $F$ the property on whose possession by the object the fittingness of instances of $E$ depends (Principle~\ref{p:fit}). That dependence is not due to an accompanying belief (Principle~\ref{p:belief}) or to a feeling (Principle~\ref{p:feel}). By Assumption~\ref{a:content}, instances of $E$ represent their objects as $F$. $F$ is evaluative: in each case it is a property in virtue of which there is reason for a response, to flee or protect oneself from what is dangerous, to avoid ingesting or touching what contaminates, to resist or protest an offence (Definition~\ref{d:eval}). And $F$ is the formal object of $E$, since an instance of $E$ that did not present its object as $F$ would not be assessable by $F$, and so would not have the fittingness conditions of the type.
\end{proof}

\begin{corollary}\label{c:fear}
Fear represents its object as dangerous, and disgust represents its object as contaminating.
\end{corollary}

The second part of the corollary agrees with the psychology of disgust, whose eliciting conditions, on the influential review of \citet{RozinFallon1987}, are organized around offensiveness and the threat of contamination by ingestion or contact.

\begin{corollary}\label{c:remorse}
Remorse represents one's own past act as having infringed, through one's own agency, a serious claim of another.
\end{corollary}

\begin{proof}
Remorse is fitting when, and only when, the act at which it is directed infringed a serious claim of another through the agent's own agency. That it is not fitting merely when the act had a bad outcome is shown by the contrast with agent-regret, which is fitting when one's act causes harm without infringing anyone's claim; that it does not require that the act was wrong is shown by cases of justified infringement, in which remorse remains fitting though guilt does not. By Theorem~\ref{t:main}, remorse represents its object as having the property on which its fittingness depends.
\end{proof}

Corollary~\ref{c:remorse} agrees with the account of remorse in ETH-27-a, on which remorse presents its object as an infringement of a claim and remains fitting in dirty-hands cases. It illustrates that the formal objects of emotions can be moral properties and need not be properties of the subject's own welfare.

The account fits the distinction between emotions and quasi-emotions in AES-05-a. On that account, a quasi-emotion of a kind shares the physiological profile of the emotion without the belief that the emotion requires, and the appreciator of a horror film experiences quasi-fear, not fear. Theorem~\ref{t:main} assigns the evaluative content to the emotion as an affective state and leaves open what further conditions individuate fear proper; the belief condition of AES-05-a is such a condition. The two accounts together say that fear represents its object as dangerous and that fear proper also involves taking the danger to be actual.

\section{Objections and replies}

\begin{objection}[The attitudinal theory]
\citet{DeonnaTeroni2015} argue that the evaluative property belongs to the attitude, not the content. Fear does not represent the dog as dangerous; it is a felt bodily stance of readiness to act so as to diminish the dog's impact, and this stance is correct when the dog is dangerous. The perceptual theory has no explanation of why the content of fear is always danger, whereas on the attitudinal theory this is a constraint of the attitude.
\end{objection}

\begin{reply}
On either view an instance of fear is correct according as its object is dangerous. The question is whether that is a fact about the state's content, and the attitudinal theory locates the evaluative property in the correctness conditions while denying that the state represents it. Assumption~\ref{a:content} says that this combination is unstable: a state correct according as its object is $F$, not in virtue of another representation, represents its object as $F$. The constraint that fear always concerns danger is explained on our view as on the attitudinal one, since the formal object of a type is fixed by the type (Definition~\ref{d:formal}). Belief is also an attitude whose correctness is truth, and that does not show that beliefs do not represent their objects. A second consideration favors content: recalcitrant emotions conflict with the subject's judgment, and the conflict is most naturally explained as a conflict between a state that presents the flight as dangerous and a belief that it is not.
\end{reply}

\begin{objection}[Objectless moods]
Anxiety, depression and elation need not be directed at anything. If they are emotions, not all emotions represent their objects as anything.
\end{objection}

\begin{reply}
Moods are not emotions in the sense of the theorem, which concerns states directed at particular objects. They are plausibly states that represent the world, or the subject's situation as a whole, as having an evaluative character: as threatening, or as without value, or as full of promise. If so, they too represent evaluative properties, with a general object. The theorem does not depend on this extension.
\end{reply}

\begin{objection}[Over-intellectualization]
Infants and animals feel fear without possessing the concept of danger. If fear represented danger, they could not feel it.
\end{objection}

\begin{reply}
The objection assumes that representing a property requires possessing a concept of it. Perceptual experience represents colors and shapes in subjects who lack the concepts of these, and the content of emotions can likewise be non-conceptual \citep{Tappolet2016}. What the infant's fear must be able to do is to be correct or incorrect according as its object is dangerous, which it can do without the infant's thinking of danger as such.
\end{reply}

\begin{objection}[The metaphysics of value]
If emotions represent evaluative properties and are often fitting, there are evaluative properties. Their nature is obscure, and a theory of emotion should not depend on a contested metaethics.
\end{objection}

\begin{reply}
The theorem requires only that evaluative predicates express properties that objects can have or lack, which most metaethical views allow. Danger and contamination are as natural as properties come; for moral and aesthetic properties, a response-dependent view, on which they are analogous to secondary qualities \citep{McDowell1985}, is sufficient. Whichever view is right, emotions represent what it says these properties are.
\end{reply}

\section{What the result does not establish}

Theorem~\ref{t:main} does not establish that emotions are perceptions, or that they justify evaluative beliefs, though it is a premise of the perceptual view and of the claim that emotions can give evidence about value. It does not establish what makes an emotion of a type fear proper as opposed to a quasi-emotion of the kind. And it does not identify the formal object of every emotion type; the cases given are those on which the literature largely agrees, and the identification of others, such as nostalgia, is a matter of further argument.

\section{Conclusion}

Emotions are correct or incorrect according as their objects have the properties specific to their types, and this correctness is neither that of the beliefs that accompany them nor that of bodily feeling. A state that is correct in this way represents its object as having the property. Emotions therefore represent evaluative properties: fear represents danger, disgust contamination, and remorse the infringement of a claim through one's own agency.

\begin{thebibliography}{99}\small
\bibitem[Brady(2009)]{Brady2009} Brady, M.~S. (2009). The irrationality of recalcitrant emotions. \emph{Philosophical Studies}, 145(3), 413--430.
\bibitem[D'Arms and Jacobson(2000)]{DArmsJacobson2000} D'Arms, J. and Jacobson, D. (2000). The moralistic fallacy: on the \enquote{appropriateness} of emotions. \emph{Philosophy and Phenomenological Research}, 61(1), 65--90.
\bibitem[de Sousa(1987)]{deSousa1987} de Sousa, R. (1987). \emph{The Rationality of Emotion}. MIT Press.
\bibitem[Deonna and Teroni(2012)]{DeonnaTeroni2012} Deonna, J.~A. and Teroni, F. (2012). \emph{The Emotions: A Philosophical Introduction}. Routledge.
\bibitem[Deonna and Teroni(2015)]{DeonnaTeroni2015} Deonna, J.~A. and Teroni, F. (2015). Emotions as attitudes. \emph{Dialectica}, 69(3), 293--311.
\bibitem[James(1884)]{James1884} James, W. (1884). What is an emotion? \emph{Mind}, 9(34), 188--205.
\bibitem[Kenny(1963)]{Kenny1963} Kenny, A. (1963). \emph{Action, Emotion and Will}. Routledge \& Kegan Paul.
\bibitem[Lazarus(1991)]{Lazarus1991} Lazarus, R.~S. (1991). \emph{Emotion and Adaptation}. Oxford University Press.
\bibitem[McDowell(1985)]{McDowell1985} McDowell, J. (1985). Values and secondary qualities. In T.~Honderich (Ed.), \emph{Morality and Objectivity}. Routledge \& Kegan Paul.
\bibitem[Nussbaum(2001)]{Nussbaum2001} Nussbaum, M.~C. (2001). \emph{Upheavals of Thought: The Intelligence of Emotions}. Cambridge University Press.
\bibitem[Prinz(2004)]{Prinz2004} Prinz, J.~J. (2004). \emph{Gut Reactions: A Perceptual Theory of Emotion}. Oxford University Press.
\bibitem[Rozin and Fallon(1987)]{RozinFallon1987} Rozin, P. and Fallon, A.~E. (1987). A perspective on disgust. \emph{Psychological Review}, 94(1), 23--41.
\bibitem[Solomon(1976)]{Solomon1976} Solomon, R.~C. (1976). \emph{The Passions}. Doubleday.
\bibitem[Tappolet(2016)]{Tappolet2016} Tappolet, C. (2016). \emph{Emotions, Values, and Agency}. Oxford University Press.
\end{thebibliography}

\end{document}
